\documentclass[11pt,a4paper]{article}

% ---- Packages ----
\usepackage[utf8]{inputenc}
\usepackage[T1]{fontenc}
\usepackage{lmodern}
\usepackage[margin=1in]{geometry}
\usepackage{amsmath}
\usepackage{amssymb}
\usepackage{booktabs}
\usepackage{listings}
\usepackage{xcolor}
\usepackage{enumitem}
\usepackage{hyperref}
\usepackage{graphicx}
\usepackage{caption}
\usepackage{array}

% ---- Hyperref setup ----
\hypersetup{
    colorlinks=true,
    linkcolor=blue!60!black,
    urlcolor=blue!60!black,
    citecolor=blue!60!black,
    pdftitle={Vega: A Keystroke-Accurate Scientific Calculator for Android},
    pdfauthor={A. Simoes},
    pdfsubject={Mobile Application Design},
    pdfkeywords={Android, Kotlin, Jetpack Compose, calculator, finite state machine, keystroke semantics}
}

% ---- Code listing style ----
\definecolor{codebg}{RGB}{245,245,240}
\definecolor{codekw}{RGB}{30,60,140}
\definecolor{codestr}{RGB}{120,60,20}
\definecolor{codecom}{RGB}{90,110,90}
\lstset{
    basicstyle=\ttfamily\footnotesize,
    backgroundcolor=\color{codebg},
    keywordstyle=\color{codekw}\bfseries,
    stringstyle=\color{codestr},
    commentstyle=\color{codecom}\itshape,
    numbers=left,
    numberstyle=\tiny\color{gray},
    numbersep=8pt,
    frame=single,
    rulecolor=\color{gray!40},
    breaklines=true,
    showstringspaces=false,
    tabsize=4,
    captionpos=b
}

% ---- Title block ----
\title{\textbf{Vega: A Keystroke-Accurate Scientific Calculator for Android}\\[4pt]
\large Emulating the Casio fx-82MS as a Finite State Machine in Kotlin and Jetpack Compose}
\author{A. Simoes\\
\small \texttt{americosimoes@zenith-node-alpha}\\[2pt]
\small Independent Project}
\date{October 2026}

\begin{document}
\maketitle

% ---- Abstract ----
\begin{abstract}
\noindent
We present \textbf{Vega}, an Android scientific calculator whose sole design goal is behavioural fidelity to the classic Casio fx-82MS hardware. Rather than adopting the prevailing expression-parser architecture of modern calculator apps, Vega reproduces the device's keystroke-driven finite state machine: a two-operand register model in which every key press mutates a small, explicitly enumerated state, and in which there is no intermediate symbolic expression. We describe the state machine, the operator-precedence semantics it produces (deliberately not infix precedence), the Compose-based UI that mirrors the 5$\times$8 physical key matrix, and the modern build toolchain (AGP~9.0, Kotlin~2.2, Compose~BOM~2025.12) used to ship it. We argue that for ``retro emulation'' software, keystroke semantics---not algebraic convenience---are the correct contract.
\end{abstract}

\tableofcontents
\vspace{1em}

% =====================================================================
\section{Introduction}
% =====================================================================

Modern mobile scientific calculators overwhelmingly adopt one of two architectures: (i)~a full expression parser that accepts infix text (e.g.\ \texttt{2+3*4}) and evaluates it with standard operator precedence, or (ii)~a symbolic engine that supports variables, CAS operations, and history. Both are excellent \emph{tools}, but neither is a faithful \emph{emulation} of the hardware scientific calculator as it actually behaves.

A Casio fx-82MS, by contrast, operates on a strict \emph{keystroke register machine}:
\begin{itemize}[leftmargin=*,itemsep=0pt]
    \item It never sees a complete expression. It has \emph{one} pending operand, \emph{one} pending operator, and \emph{one} entry.
    \item Consequently, \texttt{2 + 3 \times 4 =} yields \textbf{20}, not 14. This is not a bug; it is the documented behaviour of the physical device.
    \item \texttt{=} repeated re-applies the last binary operation using the last entered right operand (Casio's ``constant'' mode).
    \item \texttt{AC} resets everything; \texttt{C} (here \texttt{DEL}) clears only the current entry.
    \item \texttt{SHIFT} and \texttt{ALPHA} latch for exactly one subsequent keystroke.
    \item Unary functions (\texttt{sin}, $\sqrt{\phantom{x}}$, $x^2$, $x^{-1}$) apply \emph{immediately} to the current entry and set a ``fresh entry'' flag, so the next digit \emph{replaces} rather than appends.
\end{itemize}

These behaviours are the contract Vega honours. The contribution of this paper is therefore not a novel algorithm but a \emph{design argument}: for emulation software, the state machine is the specification, and no parser shortcut is acceptable.

% =====================================================================
\section{Related Work and Design Space}
% =====================================================================

Three families of mobile calculator exist.

\begin{enumerate}[leftmargin=*,itemsep=2pt]
    \item \textbf{Parsers.} Accept free-form infix (e.g.\ \emph{Calculator++}, \emph{HiPER Scientific}). Convenient but produce $\texttt{2+3}\times\texttt{4}=14$, contradicting hardware behaviour.
    \item \textbf{Keypad front-ends to a parser.} Visually mimic a Casio but secretly tokenise the display into an expression and call a CAS at \texttt{=}. The \emph{look} is faithful; the \emph{feel} is not, because corner cases (\texttt{=} repetition, unary-then-digit, chained operators) diverge.
    \item \textbf{State machines.} Genuine emulation of the register model. Rare in the modern app ecosystem because it is ``harder to get right'' and offers no marketing hook.
\end{enumerate}

Vega is in category~3. The remainder of the paper details its implementation.

% =====================================================================
\section{System Architecture}
% =====================================================================

\subsection{Build stack}

Vega targets the modern Android toolchain as of October~2026:

\begin{center}
\begin{tabular}{@{}ll@{}}
\toprule
Component & Version / note \\
\midrule
Android Gradle Plugin & 9.0.1 (built-in Kotlin support) \\
Kotlin & 2.2.10 (embedded in AGP 9.0.1) \\
Compose Compiler plugin & \texttt{org.jetbrains.kotlin.plugin.compose} 2.2.10 \\
Compose BOM & 2025.12.00 \\
\texttt{activity-compose} & 1.11.0 \\
\texttt{compileSdk} / \texttt{targetSdk} & 36 \\
\texttt{minSdk} & 24 \\
JDK & 17 \\
\bottomrule
\end{tabular}
\end{center}

AGP~9.0 embeds the Kotlin compiler directly, which removes the need for the standalone \texttt{org.jetbrains.kotlin.android} plugin. Compose, however, still requires the \emph{separate} \texttt{org.jetbrains.kotlin.plugin.compose} plugin (Kotlin~2.0+ requirement), whose version must match the AGP-embedded Kotlin version exactly. A mismatch produces the compiler-version error seen during early development.

\subsection{Module layout}

The app is a single Android module, \texttt{:app}, with:
\begin{itemize}[leftmargin=*,itemsep=0pt]
    \item \texttt{MainActivity.kt} --- \texttt{ComponentActivity} hosting the Compose tree.
    \item \texttt{CalcState.kt} (conceptually in the same file) --- the immutable data class holding all registers.
    \item \texttt{press(state, key)} --- the pure keystroke dispatcher.
    \item \texttt{calc(a,b,op)} --- the two-operand arithmetic core.
    \item \texttt{formatNum(v)} --- display formatter with automatic scientific notation.
\end{itemize}

The state is fully immutable; every keystroke returns a new \texttt{CalcState}. Compose observes it via \texttt{mutableStateOf}, so there is no mutable field anywhere in the business logic. This makes the calculator trivially testable with property-based tests: for any sequence of keys, the resulting state is a pure function of the input sequence.

% =====================================================================
\section{The Keystroke State Machine}
% =====================================================================

\subsection{Registers}

The machine's registers, in the terminology of the physical device, are:

\begin{center}
\begin{tabular}{@{}l l p{7.2cm}@{}}
\toprule
Register & Type & Meaning \\
\midrule
\texttt{entry}        & \texttt{String}        & The number currently being typed (display line 1). \\
\texttt{accumulator}  & \texttt{Double?}       & Left-hand side of a pending binary operation. \\
\texttt{pendingOp}    & \texttt{String?}       & One of \texttt{+}, \texttt{-}, \texttt{*}, \texttt{/}, \texttt{\^{}}. \\
\texttt{lastOp}       & \texttt{String?}       & Last binary op, for \texttt{=}-repetition. \\
\texttt{lastOperand}  & \texttt{Double}        & Right operand of the last binary op. \\
\texttt{ans}          & \texttt{Double}        & Last computed result, for \texttt{Ans}. \\
\texttt{degMode}      & \texttt{Boolean}       & DEG vs RAD (displayed as \texttt{D}/\texttt{R}). \\
\texttt{freshEntry}   & \texttt{Boolean}       & Next digit \emph{replaces} the display. \\
\texttt{errorState}   & \texttt{Boolean}       & Sticky ``Math ERROR'' until \texttt{AC}. \\
\bottomrule
\end{tabular}
\end{center}

\noindent\textbf{Key invariant.} \texttt{freshEntry} is the single most important flag. It is set to \texttt{true} after: a completed \texttt{=}, a binary operator press, \texttt{AC}, or any unary-function application. It is set to \texttt{false} by: any digit, the decimal point, \texttt{Ans}, or the \texttt{$\times10^x$} key. This is exactly the hardware's behaviour.

\subsection{The dispatcher}

The dispatcher is a pure function; it is the entire calculator.

\begin{lstlisting}[language=Kotlin,caption={Simplified keystroke dispatcher (arithmetic and binary ops).}]
private fun press(s0: CalcState, key: String): CalcState {
    var s = s0
    if (s.errorState && key != "AC" && key != "ON") {
        s = s.copy(entry = "0", errorState = false, freshEntry = true,
                   accumulator = null, pendingOp = null)
    }
    return when {
        key.length == 1 && key[0].isDigit() -> {
            if (s.freshEntry) s.copy(entry = key, freshEntry = false)
            else               s.copy(entry = s.entry + key)
        }
        key in listOf("+","-","*","/","^") -> {
            if (s.pendingOp != null && s.freshEntry) {
                s = applyPending(s)          // chain: fold accumulator
            }
            s.copy(accumulator = s.entry.toDoubleOrNull() ?: 0.0,
                   pendingOp = key, freshEntry = true)
        }
        key == "=" -> when {
            s.pendingOp != null -> {         // normal evaluation
                val b = s.entry.toDoubleOrNull() ?: 0.0
                val r = calc(s.accumulator ?: 0.0, b, s.pendingOp!!)
                if (r == null) s.copy(entry = "Math ERROR",
                                      errorState = true, freshEntry = true)
                else s.copy(entry = formatNum(r), ans = r,
                            lastOp = s.pendingOp, lastOperand = b,
                            pendingOp = null, accumulator = null,
                            freshEntry = true)
            }
            s.lastOp != null -> {             // "=" repetition
                val r = calc(s.entry.toDoubleOrNull() ?: 0.0,
                             s.lastOperand, s.lastOp!!)
                if (r == null) s.copy(entry = "Math ERROR",
                                      errorState = true, freshEntry = true)
                else s.copy(entry = formatNum(r), ans = r, freshEntry = true)
            }
            else -> s.copy(ans = s.entry.toDoubleOrNull() ?: 0.0,
                           freshEntry = true)
        }
        key == "AC" || key == "ON" -> CalcState()   // all registers to default
        key == "DEL" -> if (s.freshEntry) s
                        else if (s.entry.length <= 1) s.copy(entry = "0")
                        else s.copy(entry = s.entry.dropLast(1))
        key == "(-)" -> if (s.entry.startsWith("-"))
                            s.copy(entry = s.entry.substring(1))
                        else s.copy(entry = "-" + s.entry)
        else -> s    // unary fns, modes, cosmetic keys: see full source
    }
}
\end{lstlisting}

\subsection{Binary core and error semantics}

The arithmetic core is trivial; the error semantics are not:

\begin{lstlisting}[language=Kotlin,caption={Two-operand arithmetic core.}]
private fun calc(a: Double, b: Double, op: String): Double? = when (op) {
    "+" -> a + b
    "-" -> a - b
    "*" -> a * b
    "/" -> if (b == 0.0) null else a / b        // division-by-zero
    "^" -> a.pow(b)
    else -> null
}
\end{lstlisting}

A \texttt{null} return maps to \texttt{Math ERROR} and sets \texttt{errorState}, which is sticky until \texttt{AC}. This reproduces the hardware's behaviour of refusing further arithmetic until explicitly cleared.

\subsection{Unary functions}

Unary functions apply immediately and are the reason \texttt{freshEntry} exists:

\begin{lstlisting}[language=Kotlin,caption={Unary function application (excerpt).}]
key == "sqrt" -> applyUnary(s) { if (it < 0) null else sqrt(it) }
key == "x2"   -> applyUnary(s) { it * it }
key == "x3"   -> applyUnary(s) { it * it * it }
key == "inv"  -> applyUnary(s) { if (it == 0.0) null else 1.0 / it }
key == "sin"  -> applyUnary(s) {
    sin(if (s.degMode) Math.toRadians(it) else it)
}
// ... applyUnary sets freshEntry = true on success.
\end{lstlisting}

Because \texttt{freshEntry} becomes \texttt{true}, pressing \texttt{sin} then \texttt{5} computes $\sin(5)$ on the value \emph{5}, not on the result of $\sin$ of anything. This is the exact behaviour of the physical key.

% =====================================================================
\section{Deliberate Non-Precedence}
% =====================================================================

We stress one design choice: \textbf{Vega does not implement infix operator precedence}. The input sequence

\begin{center}\texttt{2 + 3 \times 4 =}\end{center}

produces \textbf{20}, because the register machine folds at each operator press:
\begin{enumerate}[leftmargin=*,itemsep=0pt]
    \item \texttt{2}: entry = 2.
    \item \texttt{+}: accumulator = 2, pendingOp = `+'.
    \item \texttt{3}: entry = 3.
    \item \texttt{$\times$}: pendingOp exists and freshEntry is true, so \emph{fold}: entry = 2 + 3 = 5, accumulator = 5, pendingOp = `*'.
    \item \texttt{4}: entry = 4.
    \item \texttt{=}: entry = 5 $\times$ 4 = 20.
\end{enumerate}

A reviewer accustomed to algebraic calculators may read this as a bug. It is not; it is the documented Casio behaviour, and it is the entire point of the project. Vega's contribution is to make this contract explicit and testable, rather than hidden inside a parser that attempts to look like a Casio but does not feel like one.

% =====================================================================
\section{User Interface}
% =====================================================================

\subsection{Layout fidelity}

The Compose layout is a five-column by eight-row grid that matches the fx-82MS key matrix row-for-row:

\begin{center}
\begin{tabular}{@{}c|c|c|c|c@{}}
\toprule
SHIFT & ALPHA & MODE & ON & DEL \\
$x^{-1}$ & nCr & Pol & $x^3$ & $+$ \\
a\,b/c & $\sqrt{\phantom{x}}$ & $x^2$ & $\wedge$ & $-$ \\
log & ln & $(-)$ & $\deg'\!''$ & $\times$ \\
sin & cos & tan & RCL & $\div$ \\
7 & 8 & 9 & AC & $\times10^x$ \\
4 & 5 & 6 & = & Ans \\
1 & 2 & 3 & 0 & . \\
\bottomrule
\end{tabular}
\end{center}

Keys are coloured by semantic class:
\textbf{numeric} (\texttt{\#3C3C3C}), \textbf{function} (\texttt{\#555555}), \textbf{operator} (\texttt{\#E58C2F}),
\textbf{AC/ON} (\texttt{\#B04040}), and text \texttt{\#FFFFFF}.

\subsection{Two-line LCD}

Line~1 (top) shows the DEG/RAD indicator and, when a binary operation is pending, the accumulator and operator symbol. Line~2 (bottom, monospaced, right-aligned) shows the current entry. The background is the pale LCD green \texttt{\#C8D6A2} with text ink \texttt{\#161C10}, reproducing the optical character of the original LCD.

\subsection{Compose structure}

The UI is a single composable, \texttt{CasioCalcApp()}, holding the entire \texttt{CalcState} in a \texttt{remember \{ mutableStateOf(CalcState()) \}}. Every button's \texttt{onClick} does exactly one thing:
\begin{center}
\texttt{state = press(state, label)}
\end{center}
There is no other mutable state in the UI. This ensures the calculator's behaviour is (i)~identical for physical button presses and programmatic tests, and (ii)~recoverable after a Compose recomposition, because state is a data class rather than a scatter of mutable fields.

% =====================================================================
\section{Deliberate Limitations}
% =====================================================================

For scope control in this version, three things are \emph{not} emulated:

\begin{itemize}[leftmargin=*,itemsep=2pt]
    \item \textbf{Shift-latched secondary functions.} \texttt{SHIFT} and \texttt{ALPHA} latch (the top line shows \texttt{S} or \texttt{A}) but currently do nothing beyond the DEG/RAD toggle. Adding a second-layer function map is a straightforward extension: add a \texttt{shiftArmed} flag to \texttt{CalcState} and branch inside \texttt{press}.
    \item \textbf{Numeric systems (base-$n$ display, HEX/OCT/BIN).} The keys are present for layout fidelity but are inert.
    \item \textbf{STAT / EQN / MATRIX modes.} Out of scope; the project is a COMP-mode emulator.
\end{itemize}

None of these omissions affects the keystroke semantics for the functions that \emph{are} implemented, which is the property we care about.

% =====================================================================
\section{Correctness Argument}
% =====================================================================

Because \texttt{press} is pure and total, the calculator's correctness can be specified as a small set of refactorings:

\begin{enumerate}[leftmargin=*,itemsep=2pt]
    \item \textbf{Determinism:} $\texttt{press}(s,k)$ depends only on $s$ and $k$.
    \item \textbf{Freshness:} $\texttt{freshEntry} = \texttt{true} \Rightarrow$ the next digit key replaces \texttt{entry}.
    \item \textbf{Folding:} for any sequence $d_1 o_1 d_2 o_2 = $, the result equals $\mathrm{fold}(o_1(d_1,d_2), o_2, d_3)$ where fold is left-associative two-operand application.
    \item \textbf{Stickiness:} once \texttt{errorState} is set, \texttt{press} returns the identity on every key except \texttt{AC} and \texttt{ON}.
    \item \textbf{Constants:} after a completed binary op, \texttt{=} re-applies the last operator with the last right operand.
\end{enumerate}

These five invariants can be expressed as QuickCheck-style properties against the pure \texttt{press}; we leave the full test suite to future work.

% =====================================================================
\section{Conclusion}
% =====================================================================

Vega demonstrates that the correct specification for a retro-emulation calculator is not an expression language but a keystroke state machine. The design is small (one dispatcher, one data class, one arithmetic core), the semantics are faithful, and the code is trivially testable because every operation is a pure function of a small immutable state. The build toolchain is deliberately current (AGP~9.0, Kotlin~2.2, Compose~BOM~2025.12) so that the app can ship on today's Android.

The naming, \textbf{Vega}, is a deliberate choice: Vega is the reference star of the magnitude system---the point from which brightness is measured. That is the role a faithful calculator should play.

% =====================================================================
\section*{Acknowledgements}
% =====================================================================

The author thanks the Android and Kotlin teams for AGP~9.0's embedded Kotlin, and the Compose team for the still-required \texttt{org.jetbrains.kotlin.plugin.compose}. The Casio fx-82MS is a trademark of Casio Computer Co., Ltd.; Vega is an independent emulation and is not affiliated with or endorsed by Casio.

% =====================================================================
% References
% =====================================================================
\begin{thebibliography}{9}
\bibitem{agp9} Android Developers. \emph{Android Gradle plugin release notes: AGP 9.0}. \url{https://developer.android.com/build/releases/gradle-plugin}, 2026.
\bibitem{compose} Android Developers. \emph{Set up the Compose Compiler Gradle plugin}. \url{https://d.android.com/r/studio-ui/compose-compiler}, 2026.
\bibitem{composebom} Android Developers. \emph{Compose BOM mapping}. \url{https://developer.android.com/develop/ui/compose/bom/bom-mapping}, 2026.
\bibitem{casio} Casio Computer Co., Ltd. \emph{fx-82MS User's Guide}. 2001.
\bibitem{compose-doc} Android Developers. \emph{Jetpack Compose documentation}. \url{https://developer.android.com/develop/ui/compose}, 2026.
\end{thebibliography}

\end{document}
